%!TEX root = ../main.tex
% Agradecimientos. Van antes del resumen, como pide la cátedra.

\begin{guia}
Breves, en primera persona y sinceros. Es la única sección del documento
donde corresponde el ``yo''. Suele agradecerse a la cátedra, al tutor o
tutora, a las personas entrevistadas o encuestadas, y al entorno personal.
Se completa al final, pero conviene dejarlo previsto desde el principio.
\end{guia}

A los profesores de la cátedra, por guiarme a lo largo de todo el
proyecto. Sus correcciones y preguntas me obligaron a justificar cada
decisión con evidencia, y este trabajo es mejor por eso.

A las personas que aceptaron ser entrevistadas, por su tiempo y por la
generosidad de proponer mejoras desde su experiencia en el aula. Y a
quienes respondieron la encuesta, que le dieron al proyecto datos reales
donde antes había supuestos.

A mi familia y a mis amigos más cercanos, que escucharon este proyecto
más veces de las que hubieran querido. Muchas de sus preguntas, hechas
con curiosidad y sin conocer los detalles técnicos, me mostraron cosas
que yo solo no veía: un supuesto que no se sostenía, un riesgo que no
había considerado o una forma más simple de explicar lo que estaba
construyendo. Gran parte de lo que Tinku es hoy se lo debo a esas
conversaciones.
